% CSE 715 — Computer Graphics
% Chapter 9: Geometric Representation — Bezier / B-Spline / Hermite Curves
% University of Chittagong, 7th Semester B.Sc. Engineering
% Source: Schaum's Outline of Computer Graphics (Xiang & Plastock), Chapter 9
%
% Compile with: pdflatex Chapter9_Solutions.tex (run twice for TOC)

\documentclass[12pt, a4paper]{article}

\usepackage[left=2.5cm, right=2.5cm, top=2.8cm, bottom=2.8cm, headheight=15pt]{geometry}
\usepackage{amsmath, amssymb}
\usepackage{tikz}
\usetikzlibrary{arrows.meta, positioning, shapes.geometric, calc, backgrounds, fit}
\usepackage{booktabs}
\usepackage{array}
\usepackage{xcolor}
\usepackage{enumitem}
\usepackage{fancyhdr}
\usepackage{titlesec}
\usepackage{mdframed}
\usepackage{tabularx}
\usepackage{hyperref}

\definecolor{sectionblue}{RGB}{10,60,110}
\definecolor{boxbg}{RGB}{242,247,255}
\definecolor{answerbg}{RGB}{240,255,240}
\definecolor{notesbg}{RGB}{255,250,230}

\mdfdefinestyle{solutionframe}{linecolor=sectionblue, backgroundcolor=boxbg, roundcorner=4pt, linewidth=1.4pt, innertopmargin=7pt, innerbottommargin=7pt, innerleftmargin=9pt, innerrightmargin=9pt}
\mdfdefinestyle{answerframe}{linecolor=green!55!black, backgroundcolor=answerbg, roundcorner=4pt, linewidth=1.4pt, innertopmargin=7pt, innerbottommargin=7pt, innerleftmargin=9pt, innerrightmargin=9pt}
\mdfdefinestyle{notesframe}{linecolor=orange!70!black, backgroundcolor=notesbg, roundcorner=4pt, linewidth=1pt, innertopmargin=5pt, innerbottommargin=5pt, innerleftmargin=8pt, innerrightmargin=8pt}
\newenvironment{solutionbox}{\begin{mdframed}[style=solutionframe]}{\end{mdframed}}
\newenvironment{answerbox}{\begin{mdframed}[style=answerframe]}{\end{mdframed}}
\newenvironment{notesbox}{\begin{mdframed}[style=notesframe]}{\end{mdframed}}

\pagestyle{fancy}
\fancyhf{}
\lhead{\small\color{sectionblue} CSE 715 --- Chapter 9: Bezier/B-Spline/Hermite}
\rhead{\small\color{sectionblue} Syllabus-Mandated Coverage}
\cfoot{\small\thepage}
\renewcommand{\headrulewidth}{0.5pt}

\titleformat{\section}{\Large\bfseries\color{sectionblue}}{}{0em}{}[\vspace{-2pt}\color{sectionblue}\rule{\linewidth}{0.8pt}]
\titleformat{\subsection}{\large\bfseries\color{sectionblue!85!black}}{}{0em}{}

\begin{document}

\begin{center}
  {\LARGE\bfseries\color{sectionblue} CSE 715: Computer Graphics}\\[0.6em]
  {\Large\bfseries Chapter 9 --- Geometric Representation: Bezier, B-Spline, Hermite Curves}\\[0.3em]
  {\large Syllabus-Mandated Theory + Solved Exercises}\\[0.2em]
  {\normalsize University of Chittagong \quad|\quad 7\textsuperscript{th} Semester B.Sc.\ (Engg.) Examination}\\[0.3em]
  {\small Source: Schaum's Outline of Computer Graphics (Xiang \& Plastock), Chapter 9 --- notation follows the book exactly ($H_i(x)$, $B_{i,n}(x)$, $BE_{k,n}(x)$, $P(t)$)}
\end{center}

\vspace{0.3em}\noindent\rule{\linewidth}{1pt}\vspace{0.5em}

\begin{notesbox}
\textbf{Why this document looks different from the other chapter solutions:} Bezier/B-spline/Hermite curves have \textbf{0/5 hits} in the 2020--2024 past papers --- confirmed by direct re-verification of every question in every paper this exam season. \texttt{Syllabus.txt} explicitly overrides the "skip if not in past papers" rule for this one topic: \textit{"Chapter 9: Bezier Curve, solve problems, B-spline, Hermite Curve, exercise and solved questions (other topics only previous years questions if any)."} So instead of a past-paper year-by-year breakdown, this document gives the book's core definitions plus worked textbook-style exercises, matching what the syllabus asks for. \textbf{Given the exam is imminent and this is a 0-yield topic, treat this as a stretch-goal skim, not a priority revision target} --- Z-Buffer, Cohen-Sutherland, and the other TIER 1 topics come first.
\end{notesbox}

\tableofcontents
\newpage

%=====================================================================
\section{Core Definitions}
%=====================================================================

%----------------------------------------------------------------------
\subsection{Hermite Cubic Polynomials}
%----------------------------------------------------------------------
\begin{answerbox}
Given knots $t_{i-1},t_i,t_{i+1}$, the Hermite cubic basis functions are
$$H_i(x) = \begin{cases}-\dfrac{2(x-t_{i-1})^3}{(t_i-t_{i-1})^3}+\dfrac{3(x-t_{i-1})^2}{(t_i-t_{i-1})^2} & t_{i-1}\le x\le t_i\\[6pt] -\dfrac{2(t_{i+1}-x)^3}{(t_{i+1}-t_i)^3}+\dfrac{3(t_{i+1}-x)^2}{(t_{i+1}-t_i)^2} & t_i\le x\le t_{i+1}\end{cases} \qquad \bar H_i(x) = \begin{cases}\dfrac{(x-t_{i-1})^2(x-t_i)}{(t_i-t_{i-1})^2} & t_{i-1}\le x\le t_i\\[6pt] \dfrac{(x-t_i)(t_{i+1}-x)^2}{(t_{i+1}-t_i)^2} & t_i\le x\le t_{i+1}\end{cases}$$
\textbf{Hermitian cubic interpolation:} given data points $(x_0,y_0),\dots,(x_n,y_n)$ and prescribed derivative (tangent-slope) values $y_0',\dots,y_n'$ at each point, with knots $t_i=x_i$:
$$H(x) = \sum_{i=0}^{n}\left[y_i H_i(x) + y_i'\bar H_i(x)\right]$$
This constructs a piecewise-cubic curve that passes exactly through every data point \emph{and} has the prescribed tangent slope there.
\end{answerbox}

%----------------------------------------------------------------------
\subsection{B-Splines}
%----------------------------------------------------------------------
\begin{answerbox}
For a knot set $t_0,t_1,t_2,\dots$, the $n$th-degree B-splines $B_{i,n}$ are defined \textbf{recursively}:
$$B_{i,0}(x) = \begin{cases}1 & t_i\le x\le t_{i+1}\\ 0 & \text{otherwise}\end{cases} \qquad B_{i,n}(x) = \frac{x-t_i}{t_{i+n}-t_i}B_{i,n-1}(x) + \frac{t_{i+n+1}-x}{t_{i+n+1}-t_{i+1}}B_{i+1,n-1}(x)$$
$B_{i,n}(x)$ is nonzero only on $[t_i,t_{i+n+1}]$ (\textbf{local support} --- the key advantage over Bernstein/Bezier). For nonrepeated knots, $B_{i,n}(t_i)=0=B_{i,n}(t_{i+n+1})$ for $n\ge1$. Repeated knots are allowed (using the convention $\tfrac00=0$).
\textbf{Spline of degree $m$:} an $m$th-degree spline through $n+1$ data points is $S_m(x)=\sum_{i=0}^{m+n-1}a_iB_{i,m}(x)$. For the \textbf{cubic spline} ($m=3$): $n+3$ coefficients $a_i$ need $n+3$ equations --- $n+1$ come from the interpolation criterion $S_3(x_j)=y_j$, and the remaining 2 from a chosen \textbf{boundary condition}: natural ($S_3''(x_0)=S_3''(x_n)=0$), clamped (prescribed end slopes), cyclic (for closed curves), or anticyclic.
\end{answerbox}

%----------------------------------------------------------------------
\subsection{Bernstein Polynomials}
%----------------------------------------------------------------------
\begin{answerbox}
Degree-$n$ Bernstein polynomials over $[0,1]$:
$$BE_{k,n}(x) = \frac{n!}{k!(n-k)!}x^k(1-x)^{n-k}, \qquad 0\le x\le1$$
\textbf{Cubic ($n=3$) Bernstein basis:} $BE_{0,3}=1-3x+3x^2-x^3,\quad BE_{1,3}=3(x-2x^2+x^3),\quad BE_{2,3}=3(x^2-x^3),\quad BE_{3,3}=x^3$.
\end{answerbox}

%----------------------------------------------------------------------
\subsection{Bezier-Bernstein Approximation (the "Bezier Curve")}
%----------------------------------------------------------------------
\begin{answerbox}
Given a guiding polyline (control points) $P_0,\dots,P_n$, the \textbf{Bezier curve} is
$$P(t): \begin{cases}x(t)=\sum_{i=0}^{n}x_iBE_{i,n}(t)\\ y(t)=\sum_{i=0}^{n}y_iBE_{i,n}(t)\\ z(t)=\sum_{i=0}^{n}z_iBE_{i,n}(t)\end{cases} \qquad 0\le t\le1$$
\textbf{Four properties (frequently asked as "explain properties of Bezier curve"):}
\begin{enumerate}[noitemsep]
  \item \textbf{Endpoint interpolation:} $P_0=P(0),\ P_n=P(1)$ --- the curve touches the first and last control points exactly (but not the interior ones, in general).
  \item \textbf{Tangent direction at the endpoints} matches the guiding polyline's first/last segment: $P'(0)=n\cdot(P_1-P_0)$, $P'(1)=n\cdot(P_n-P_{n-1})$.
  \item \textbf{Convex hull property:} the entire Bezier curve lies within the convex hull (the "rubber band" polygon) of the control points $P_0,\dots,P_n$.
  \item \textbf{Suited to interactive/piecewise design:} multiple Bezier curves can be joined with continuous tangents by keeping the two adjacent guiding-polyline edges at the shared endpoint collinear.
\end{enumerate}
\end{answerbox}

%----------------------------------------------------------------------
\subsection{Bezier-B-Spline Approximation}
%----------------------------------------------------------------------
\begin{answerbox}
Replaces Bernstein polynomials with B-splines for \textbf{local control}:
$$P(t): \begin{cases}x(t)=\sum_{i=0}^{n}x_iB_{i,m}(t)\\ y(t)=\sum_{i=0}^{n}y_iB_{i,m}(t)\\ z(t)=\sum_{i=0}^{n}z_iB_{i,m}(t)\end{cases} \qquad 0\le t\le n-m+1$$
\textbf{Five properties:}
\begin{enumerate}[noitemsep]
  \item Shares all Bezier-Bernstein properties (endpoint interpolation, matching end-tangents, convex hull) --- in fact identical to it when $m=n$.
  \item If $m+1$ consecutive control points are collinear, that span of the curve is linear.
  \item \textbf{Local control:} moving one control point only changes $m+1$ nearby spans, not the whole curve (unlike Bezier-Bernstein, where moving any control point can reshape the entire curve).
  \item Fits the guiding polygon more closely than Bezier-Bernstein.
  \item Control points can be given \textbf{multiplicity} $k\ge2$ --- with multiplicity $m+1$, the curve is pulled all the way through that control point.
\end{enumerate}
\textbf{Cavalier/cabinet-style memory hook:} Bezier-Bernstein = "one control point moves, whole curve reshapes" (global); Bezier-B-spline = "one control point moves, only nearby span reshapes" (local) --- this local-vs-global contrast is the single most quotable fact if only one property can be remembered.
\end{answerbox}

%=====================================================================
\section{Solved Exercises (adapted from Schaum's Solved Problems 9.3, 9.4, 9.6)}
%=====================================================================

\subsection*{Exercise 1 --- B-spline is Zero Outside Its Support Interval}
\begin{solutionbox}
\textbf{Show that} $B_{i,n}(x)=0$ if $x<t_i$ or $x>t_{i+n+1}$.
\end{solutionbox}
\begin{answerbox}
By induction on the degree $n$, base case $n=0$: illustrate the first step for degree 1, $B_{i,1}(x)$. Suppose $x<t_i$; then also $x<t_{i+1}$, so both zero-degree splines $B_{i,0}(x)=0$ and $B_{i+1,0}(x)=0$ (since $x$ is left of both their support intervals). By the recursive formula, $B_{i,1}(x)$ is a combination of $B_{i,0}(x)$ and $B_{i+1,0}(x)$, both zero, so $B_{i,1}(x)=0$.
For the boundary case $x=t_i$ (nonrepeated knots): $B_{i,0}(t_i)=1$ but $B_{i+1,0}(t_i)=0$, giving
$$B_{i,1}(t_i) = \frac{t_i-t_i}{t_{i+1}-t_i}(1) + \frac{t_{i+2}-t_i}{t_{i+2}-t_{i+1}}(0) = 0$$
Similar reasoning handles $x\ge t_{i+2}$. The same induction pattern extends to any degree $n$: each step only ever combines two lower-degree splines that are individually zero outside their own (narrower) support, so the zero-outside-support property propagates upward.
\end{answerbox}

\subsection*{Exercise 2 --- Knot Set for Cubic B-Spline Interpolation}
\begin{solutionbox}
\textbf{Given} data points $P_0(0,0),P_1(1,2),P_2(2,1),P_3(3,-1),P_4(4,10),P_5(5,5)$, find a knot set $t_0,\dots,t_9$ for cubic B-spline interpolation ($m=3,\ n=5$).
\end{solutionbox}
\begin{answerbox}
Need $m+n+2=10$ knots $t_0,\dots,t_9$. \textbf{Scheme:} repeat the boundary knots $m+1=4$ times at each end:
$$t_0=t_1=t_2=t_3=-1\ (<x_0=0) \qquad t_6=t_7=t_8=t_9=6\ (>x_5=5)$$
Remaining interior knots via $t_{i+4}=\dfrac{x_{i+1}+x_{i+2}+x_{i+3}}{3}$ for $i=0,1$ (the averaging scheme) or the simpler alternative $t_{i+4}=x_{i+2}$:
$$t_4 = x_2 = 2, \qquad t_5 = x_3 = 3$$
\textbf{Final knot set:} $(-1,-1,-1,-1,2,3,6,6,6,6)$.
\end{answerbox}

\subsection*{Exercise 3 --- Bezier Curve: Endpoint and Midpoint Evaluation}
\begin{notesbox}
Illustrative exercise (not from a past paper) showing the standard "evaluate a cubic Bezier curve at a given $t$" computation --- a natural professor question given the syllabus explicitly asks for "solve problems" on Bezier curves.
\end{notesbox}
\begin{solutionbox}
\textbf{Given} cubic Bezier control points $P_0(0,0),\ P_1(1,3),\ P_2(3,3),\ P_3(4,0)$, find the curve point at $t=0.5$.
\end{solutionbox}
\begin{answerbox}
Using the cubic Bernstein basis from \S1.3: at $t=0.5$, $BE_{0,3}=1-3(.5)+3(.25)-.125=0.125$, $BE_{1,3}=3(.5-2(.25)+.125)=3(.125)=0.375$, $BE_{2,3}=3(.25-.125)=3(.125)=0.375$, $BE_{3,3}=.125$. (Check: sum $=0.125+0.375+0.375+0.125=1$ \checkmark.)
$$x(0.5) = 0(.125)+1(.375)+3(.375)+4(.125) = 0+0.375+1.125+0.5 = 2.0$$
$$y(0.5) = 0(.125)+3(.375)+3(.375)+0(.125) = 1.125+1.125 = 2.25$$
So $P(0.5)=(2.0,\,2.25)$. \textbf{Sanity checks from the properties in \S1.4:} $P(0)=P_0=(0,0)$ \checkmark (endpoint interpolation); the curve point at $t=0.5$ lies inside the convex hull (the quadrilateral $P_0P_1P_2P_3$) \checkmark.
\end{answerbox}

%=====================================================================
\section{Quick Summary}
%=====================================================================
\begin{center}
\renewcommand{\arraystretch}{1.3}
\begin{tabular}{p{4.2cm}|p{9.3cm}}
\toprule
\textbf{Concept} & \textbf{One-line memory hook} \\
\midrule
Hermite cubic & Interpolates through points \emph{and} matches prescribed tangent slopes there \\
B-spline (general) & Recursive basis, built up from degree 0; \textbf{local support} on $[t_i,t_{i+n+1}]$ \\
Bernstein polynomial & $BE_{k,n}(x)=\binom{n}{k}x^k(1-x)^{n-k}$ --- the basis functions behind Bezier \\
Bezier-Bernstein curve & Endpoints interpolated, tangents match polyline ends, convex hull, but \textbf{global} control \\
Bezier-B-spline curve & Same nice properties as Bezier, but \textbf{local} control + closer polygon fit + control-point multiplicity \\
\bottomrule
\end{tabular}
\end{center}
\begin{notesbox}
\textbf{Exam-day reality check:} this topic has never appeared in 5 years of past papers. If it's the night before the exam and time is short, skip straight back to Z-Buffer, Cohen-Sutherland, Cavalier/Cabinet, and the other TIER 1/TIER 2 topics --- come back here only if all of those are solid.
\end{notesbox}

\end{document}
